% Generated draft: requires mathematical and source review.
\documentclass[11pt]{article}
\usepackage[margin=25mm]{geometry}
\usepackage{fontspec}
\setmainfont{Latin Modern Roman}
\usepackage{amsmath,amssymb,mathtools,unicode-math}
\setmathfont{Latin Modern Math}
\usepackage{enumitem,needspace,xurl,hyperref}
\setlength{\parindent}{0pt}
\setlength{\parskip}{5pt}
\setlength{\emergencystretch}{3em}
\newcommand{\N}{\mathbb{N}}
\newcommand{\Z}{\mathbb{Z}}
\newcommand{\Q}{\mathbb{Q}}
\newcommand{\R}{\mathbb{R}}
\newcommand{\C}{\mathbb{C}}
\newcommand{\F}{\mathbb{F}}
\newcommand{\A}{\mathbb{A}}
\newcommand{\PP}{\mathbb{P}}
\newcommand{\E}{\mathbb{E}}
\newcommand{\eps}{\varepsilon}
\newcommand{\dd}{\,\mathrm{d}}
\begin{document}
\section*{UnsolvedMath ID: 30006515}
\textbf{Source problem code:} catalogue entry
\subsection{Definitions and mathematical statement}
Let \(A\) be an infinite first-order structure with underlying set \(|A|\). Define
\[
\operatorname{Aut}(A)=\{g:|A|\to |A|:g\text{ is an automorphism of }A\}
\]
and
\[
\operatorname{Emb}_{\mathrm{el}}(A)=\{e:|A|\to |A|:e\text{ is an elementary embedding}\}.
\]
Here an elementary embedding is an injective map such that, for every first-order formula \(\varphi(x_1,\ldots,x_n)\) in the language of \(A\) and every tuple \(\bar a\in |A|^n\),
\[
A\models\varphi(\bar a)\quad\Longleftrightarrow\quad A\models\varphi(e(\bar a)).
\]
Composition makes these objects a group and a monoid, respectively. For a conventional precise interpretation of their topology, give \(|A|\) the discrete topology and give both objects the subspace topology inherited from \(|A|^{|A|}\). Thus basic neighborhoods of a map \(f\) specify its values on a finite subset \(F\subseteq |A|\):
\[
U(f,F)=\{h:h\!\upharpoonright_F=f\!\upharpoonright_F\},
\]
where \(h\) ranges over the relevant group or monoid. The supplied question does not explicitly specify this topology, so its use here is a convention requiring confirmation.

\textbf{Reconstruction question.} Determine what information about \(A\) is determined by the topological group \(\operatorname{Aut}(A)\), and separately by the topological monoid \(\operatorname{Emb}_{\mathrm{el}}(A)\). More precisely, for a specified class \(\mathcal K\) of infinite structures, determine the consequences for \(A,B\in\mathcal K\) of each of
\[
\operatorname{Aut}(A)\cong_{\mathrm{top.grp}}\operatorname{Aut}(B),
\qquad
\operatorname{Emb}_{\mathrm{el}}(A)\cong_{\mathrm{top.mon}}\operatorname{Emb}_{\mathrm{el}}(B).
\]
These denote algebraic isomorphisms that are homeomorphisms. Equivalently, determine which properties \(P\) satisfy, for all \(A,B\in\mathcal K\),
\[
\operatorname{Aut}(A)\cong_{\mathrm{top.grp}}\operatorname{Aut}(B)
\ \Longrightarrow\ (P(A)\Longleftrightarrow P(B)),
\]
and which satisfy the analogous implication for elementary-embedding monoids.

The accompanying talk title specifically concerns \(\omega\)-categorical structures: a structure is \(\omega\)-categorical when its complete first-order theory has exactly one countably infinite model up to isomorphism. The broader supplied statement concerns infinite structures without this restriction. The intended class \(\mathcal K\), language restrictions, cardinalities, and precise reconstruction target are unresolved. In particular, the record does not specify a single conjectured implication to isomorphism or to bi-interpretability. Bi-interpretability means mutual first-order interpretability with the composite interpretations definably isomorphic to the respective original structures; the permitted use of parameters is also unspecified.

The supplied dated curation reports reconstruction up to bi-interpretability in the \(\omega\)-categorical setting and says that stronger exact reconstruction is not general. This is recorded as a qualified source assessment, not asserted here as a verified theorem or as a fully specified conjecture.
\subsection{Short English statement}
Ask which features of an infinite structure are determined by its automorphism group or its monoid of elementary self-embeddings, including their topology. The talk context singles out \(\omega\)-categorical structures, but the supplied question leaves the comparison class and the desired notion of reconstruction unspecified. This draft therefore states a reconstruction programme requiring human mathematical review.
\subsection{Sources}
\begin{itemize}
\item Local UnsolvedMath registry, supplied complete source record, permanent ID \(30006515\), source problem code \texttt{catalogue entry}, title: \emph{Reconstructing Structures from Topological Symmetry Monoids}. The registry labels the problem \texttt{partially\_solved} and the statement \texttt{corrected\_verified}; these are supplied curation labels. Its statement-verification field reports that the official report was checked and the question isolated and repaired; no independent source verification is claimed here.
\item Supplied report citation: \emph{Homogeneous Structures: Model Theory meets Universal Algebra (2026)}, Oberwolfach Reports, pp.\ 3063--3124, DOI \texttt{10.4171/owr/2025/58}. \url{https://doi.org/10.4171/owr/2025/58}. The registry separately gives proposed year 2025; the relationship between these dates is not resolved in this draft.
\item Supplied original-statement attribution: Paolo Marimon, \emph{On the topological reconstruction of monoids of elementary embeddings of \(\omega\)-categorical structures}, joint work with J.\ de la Nuez Gonzales, Zaniar Ghadernezhad, and Michael Pinsker. The registry has no separate named proposer.
\item Supplied literature reference: Ahlbrandt--Ziegler, \emph{Quasi finitely axiomatizable totally categorical theories}. \url{https://doi.org/10.1016/0168-0072(86)90037-0}.
\item Supplied literature reference: de la Nuez Gonzalez et al., \emph{Minimal and intrinsic topologies on monoids}. \url{https://arxiv.org/abs/2603.28419}.
\item Literature assessment dated exactly \texttt{2026-08-22}: the supplied curation labels the problem partially solved and states that topological automorphism groups and elementary-embedding monoids recover \(\omega\)-categorical structures up to bi-interpretability, while stronger exact reconstruction is not general. This is supplied dated curation derived from the OWR record, not an independently verified assessment. Its time-sensitive status and the hypotheses and scope of the cited claims require independent checking.
\end{itemize}
\end{document}
